\paragrafos{Συλλογή συνεχών ποσοτικών δεδομένων}{Γ.1.2, «Συλλογή συνεχών ποσοτικών δεδομένων»}

\begin{anagke}
Ο Ηλίας φέρνει στην τάξη μια μεζούρα και μετρά το ύψος και των $25$ παιδιών της Α$3$, σε εκατοστά. Γράφει τις μετρήσεις στο τετράδιό του, με τη σειρά που στάθηκαν τα παιδιά στην ουρά:

\noindent\mbox{}\hfill 160 \ 165 \ 154 \ 169 \ 171 \ 162 \ 155 \ 174 \ 157 \ 154 \ 161 \ 148 \ 163\hfill\mbox{}\\
\mbox{}\hfill 156 \ 160 \ 167 \ 158 \ 152 \ 163 \ 151 \ 150 \ 166 \ 159 \ 153 \ 158\hfill\mbox{}

«Αυτό δεν διαβάζεται!» λέει η Ρήνη. «Σχεδόν κάθε παιδί έχει διαφορετικό ύψος. Αν φτιάξω πίνακα με μια γραμμή για κάθε τιμή, θα έχω $21$ γραμμές με συχνότητα $1$ ή $2$». Η Φανή ρωτά: «Δεν μπορούμε να βάλουμε μαζί τα παιδιά με ύψος κοντά το ένα στο άλλο;».

Σε αυτή την παράγραφο μαθαίνουμε να ομαδοποιούμε τα συνεχή δεδομένα σε κλάσεις και να φτιάχνουμε πίνακα συχνοτήτων και σχετικών συχνοτήτων που διαβάζεται με μια ματιά.
\end{anagke}

\begin{lexi}{κλάση, ομαδοποίηση, κέντρο, πλάτος}{ονόματα για τις ομάδες}
Η \lexiword{κλάση} σημαίνει την «τάξη», την ομάδα: όπως η τάξη του σχολείου μαζεύει παιδιά ίδιας ηλικίας, η κλάση μαζεύει τιμές που είναι κοντά. Η \lexiword{ομαδοποίηση} είναι το «βάζω σε ομάδες». Το \lexiword{κέντρο} της κλάσης είναι το μέσο της και το \lexiword{πλάτος} της είναι πόσο φαρδιά είναι.
\end{lexi}

\begin{metafrasi}
Βάζω σε ομάδες τιμές που είναι κοντά. & ομαδοποίηση, κλάσεις\\
Η ομάδα από $150$ μέχρι $155$ (χωρίς το $155$). & κλάση $150-155$\\
Πόσο φαρδιά είναι η κλάση. & πλάτος κλάσης\\
Το μέσο της κλάσης. & κέντρο κλάσης\\
Πόσα στοιχεία έχει η κλάση. & συχνότητα της κλάσης\\
Ποιο μέρος του συνόλου είναι η κλάση. & σχετική συχνότητα της κλάσης\\
\end{metafrasi}

\tmima{Η θεωρία}

Στα συνεχή ποσοτικά δεδομένα (ύψος, βάρος, χρόνος) οι διαφορετικές τιμές είναι πάρα πολλές και σχεδόν δεν επαναλαμβάνονται. Τότε είναι προτιμότερο να χωρίσουμε τις παρατηρήσεις σε ομάδες που λέγονται \textbf{κλάσεις}. Η διαδικασία αυτή ονομάζεται \textbf{ομαδοποίηση παρατηρήσεων}.

\begin{orismos}{κλάσεις ίσου πλάτους}
Κάθε κλάση περιγράφεται από τα δύο άκρα της, π.χ. $150-155$. \textbf{Πλάτος} της κλάσης είναι η διαφορά του μικρότερου άκρου από το μεγαλύτερο: $155-150=5$. \textbf{Κέντρο} της κλάσης είναι ο αριθμός που βρίσκεται στο μέσο της: $\frac{150+155}{2}=152{,}5$. Οι κλάσεις που φτιάχνουμε έχουν \emph{ίσο πλάτος}, δεν επικαλύπτονται και καλύπτουν όλες τις παρατηρήσεις.
\end{orismos}

\begin{kanonas}{πού μπαίνει μια τιμή}
Μια παρατήρηση που συμπίπτει με το \emph{δεξί} άκρο μιας κλάσης τοποθετείται στην \emph{επόμενη} κλάση. Για παράδειγμα, η τιμή $155$ ανήκει στην κλάση $155-160$ και όχι στην $150-155$.
\end{kanonas}

\begin{kanonas}{πώς φτιάχνουμε τις κλάσεις}
$1.$ Βρίσκουμε τη μικρότερη και τη μεγαλύτερη τιμή. $2.$ Διαλέγουμε ένα πλάτος που δίνει περίπου $5$ έως $7$ κλάσεις και αριθμούς εύκολους (π.χ. $5$ ή $10$). $3.$ Ξεκινάμε από έναν στρογγυλό αριθμό λίγο μικρότερο ή ίσο με τη μικρότερη τιμή και προχωρούμε με το πλάτος ώσπου να καλύψουμε και τη μεγαλύτερη. $4.$ Μετράμε πόσες τιμές πέφτουν σε κάθε κλάση.
\end{kanonas}

\begin{kanonas}{σχετική συχνότητα}
$$\text{σχετική συχνότητα της κλάσης}=\frac{\text{συχνότητα της κλάσης}}{\text{σύνολο παρατηρήσεων}}$$
Το αποτέλεσμα γράφεται ως δεκαδικός ή, συνήθως, ως ποσοστό: $\frac{6}{25}=0{,}24=24\%$.
\end{kanonas}

Ο πίνακας με τις κλάσεις, τις συχνότητες και τις σχετικές συχνότητες λέγεται \textbf{πίνακας κατανομής συχνοτήτων}. Στην τελευταία γραμμή γράφουμε το σύνολο: το άθροισμα των συχνοτήτων είναι το πλήθος των παρατηρήσεων και το άθροισμα των σχετικών συχνοτήτων είναι $100\%$.

\begin{prosoxi}
\begin{itemize}[leftmargin=1.2em,itemsep=0.3ex,topsep=2pt]
\item Η συχνότητα και η σχετική συχνότητα ενός πίνακα με κλάσεις αναφέρονται \emph{στην κλάση} και όχι σε συγκεκριμένους αριθμούς. Αν η κλάση $165-170$ έχει συχνότητα $4$, δεν ξέρουμε πόσα παιδιά έχουν ύψος ακριβώς $167$.
\item Το δεξί άκρο δεν ανήκει στην κλάση. Το $160$ μπαίνει στην $160-165$.
\item Οι κλάσεις πρέπει να έχουν το ίδιο πλάτος και να μην αφήνουν κενά ούτε να επικαλύπτονται.
\item Το άθροισμα των σχετικών συχνοτήτων πρέπει να βγαίνει $100\%$. Αν δεν βγαίνει, έχει γίνει λάθος στην καταμέτρηση ή στη διαίρεση.
\end{itemize}
\end{prosoxi}

\tmima{Λυμένα παραδείγματα}

\paradeigma{Τα ύψη της Α$3$}
\tixero Οι $25$ μετρήσεις του Ηλία. Η μικρότερη είναι $148$ και η μεγαλύτερη $174$.\par
\tithelo Πίνακας κατανομής συχνοτήτων με κλάσεις πλάτους $5$ cm.\par
\possindeo Οι τιμές είναι από $148$ έως $174$, άρα οι κλάσεις ξεκινούν από το $145$ και φτάνουν ως το $175$: έξι κλάσεις. Μετράμε κάθε κλάση. Η $150-155$ περιέχει τα $150$, $151$, $152$, $153$, $154$, $154$: έξι τιμές. Το $155$ δεν είναι εδώ, ανήκει στην επόμενη κλάση.\par
\apantisi

\begin{center}
\renewcommand{\arraystretch}{1.25}
\begin{tabular}{cccc}
\toprule
\textbf{Κλάσεις} & \textbf{Κέντρο} & \textbf{Συχνότητα} & \textbf{Σχετική συχνότητα}\\
\midrule
$145-150$ & $147{,}5$ & $1$ & $4\%$\\
$150-155$ & $152{,}5$ & $6$ & $24\%$\\
$155-160$ & $157{,}5$ & $6$ & $24\%$\\
$160-165$ & $162{,}5$ & $6$ & $24\%$\\
$165-170$ & $167{,}5$ & $4$ & $16\%$\\
$170-175$ & $172{,}5$ & $2$ & $8\%$\\
\midrule
\textbf{Σύνολο} & & $25$ & $100\%$\\
\bottomrule
\end{tabular}
\end{center}

\elegxos $1+6+6+6+4+2=25$ και $4+24+24+24+16+8=100$.

\paradeigma{Ο χρόνος διαδρομής}
\tixero Η Φανή ρώτησε $20$ παιδιά πόσα λεπτά κάνουν για να φτάσουν στο σχολείο:\par
\noindent\mbox{}\hfill 15 \ 8 \ 14 \ 25 \ 10 \ 18 \ 11 \ 16 \ 12 \ 9 \ 7 \ 5 \ 10 \ 3 \ 6 \ 20 \ 13 \ 4 \ 22 \ 9\hfill\mbox{}\par
\tithelo Πίνακας με $6$ κλάσεις ίσου πλάτους.\par
\possindeo Οι τιμές είναι από $3$ έως $25$. Διαλέγουμε πλάτος $5$ και ξεκινάμε από το $0$: $0-5$, $5-10$, $10-15$, $15-20$, $20-25$, $25-30$. Η τιμή $5$ πάει στην $5-10$, οι $10$ στην $10-15$, η $15$ στην $15-20$, η $20$ στην $20-25$ και η $25$ στην $25-30$. Σε κάθε σύνολο μετράμε: $\{3,4\}$, $\{5,6,7,8,9,9\}$, $\{10,10,11,12,13,14\}$, $\{15,16,18\}$, $\{20,22\}$, $\{25\}$.\par
\apantisi

\begin{center}
\renewcommand{\arraystretch}{1.25}
\begin{tabular}{cccc}
\toprule
\textbf{Κλάσεις (λεπτά)} & \textbf{Κέντρο} & \textbf{Συχνότητα} & \textbf{Σχετική συχνότητα}\\
\midrule
$0-5$ & $2{,}5$ & $2$ & $10\%$\\
$5-10$ & $7{,}5$ & $6$ & $30\%$\\
$10-15$ & $12{,}5$ & $6$ & $30\%$\\
$15-20$ & $17{,}5$ & $3$ & $15\%$\\
$20-25$ & $22{,}5$ & $2$ & $10\%$\\
$25-30$ & $27{,}5$ & $1$ & $5\%$\\
\midrule
\textbf{Σύνολο} & & $20$ & $100\%$\\
\bottomrule
\end{tabular}
\end{center}

\elegxos $2+6+6+3+2+1=20$ και $10+30+30+15+10+5=100$.

\paradeigma{Διαβάζω τον πίνακα}
\tixero Ο πίνακας του προηγούμενου παραδείγματος.\par
\tithelo Πόσα παιδιά κάνουν λιγότερο από $10$ λεπτά; Ποιο ποσοστό κάνει τουλάχιστον $15$ λεπτά;\par
\possindeo «Λιγότερο από $10$ λεπτά» σημαίνει οι κλάσεις $0-5$ και $5-10$: $2+6=8$ παιδιά. «Τουλάχιστον $15$» σημαίνει οι κλάσεις $15-20$, $20-25$, $25-30$: $3+2+1=6$ παιδιά, δηλαδή $\frac{6}{20}=0{,}3=30\%$. Ή, με τις σχετικές συχνότητες: $15\%+10\%+5\%=30\%$.\par
\apantisi $8$ παιδιά και $30\%$.\par
\elegxos Τα παιδιά που κάνουν λιγότερο από $10$ λεπτά είναι $40\%$ του συνόλου και τα υπόλοιπα $60\%$.

\paradeigma{Οι κλάσεις μιας μέτρησης}
\tixero Οι μετρήσεις είναι από $12$ έως $47$ και θέλουμε κλάσεις πλάτους $10$.\par
\tithelo Οι κλάσεις και το κέντρο καθεμιάς.\par
\possindeo Ξεκινάμε από το $10$: $10-20$, $20-30$, $30-40$, $40-50$. Το $47$ είναι μέσα στην τελευταία. Το κέντρο κάθε κλάσης είναι το μέσο των άκρων της.\par
\apantisi Τέσσερις κλάσεις με κέντρα $15$, $25$, $35$, $45$.\par
\elegxos Κάθε κέντρο απέχει $5$ (το μισό του πλάτους) από τα άκρα της κλάσης του.

\tmima{Ασκήσεις}

\zoni{Το κατάλαβα}{κλάσεις, κέντρα, συχνότητες}

\ex Σε ποια κλάση από τις $0-10$, $10-20$, $20-30$, $30-40$ ανήκει κάθε τιμή: α) $9$, β) $10$, γ) $19$, δ) $20$, ε) $3$, στ) $30$, ζ) $38$, η) $29{,}9$;

\ex Βρες το πλάτος και το κέντρο των κλάσεων: α) $30-40$, β) $150-160$, γ) $1-1{,}5$, δ) $0-25$.

\ex Συμπλήρωσε τον πίνακα, όπου η μια στήλη έχει μόνο τις συχνότητες $3$, $7$, $6$, $4$ για τις κλάσεις $0-5$, $5-10$, $10-15$, $15-20$: το σύνολο και τις σχετικές συχνότητες.

\ex Η μικρότερη τιμή είναι $7$ και η μεγαλύτερη $28$. Ποιες κλάσεις πλάτους $5$ καλύπτουν όλα τα δεδομένα, αν ξεκινήσουμε από το $5$; Πόσες είναι;

\ex Σε σύνολο $40$ παρατηρήσεων οι συχνότητες πέντε κλάσεων είναι $4$, $10$, $16$, $8$, $2$. Βρες τις σχετικές συχνότητες ως ποσοστά.

\ex Σε σύνολο $50$ παρατηρήσεων οι σχετικές συχνότητες πέντε κλάσεων είναι $8\%$, $20\%$, $36\%$, $24\%$, $12\%$. Βρες τις συχνότητες.

\ex Σωστό ή λάθος; α) Η τιμή $160$ ανήκει στην κλάση $160-165$. β) Το κέντρο της κλάσης $10-20$ είναι $15$. γ) Το άθροισμα των σχετικών συχνοτήτων είναι πάντα $100\%$. δ) Από έναν πίνακα με κλάσεις μπορούμε να βρούμε την ακριβή τιμή κάθε παρατήρησης. ε) Αν μια κλάση έχει συχνότητα $6$ σε σύνολο $24$, η σχετική της συχνότητα είναι $25\%$.

\zoni{Το χρησιμοποιώ}{ομαδοποιώ δικά μου δεδομένα}

\ex \textbf{Το κινητό.} Ο Στέλιος ρώτησε $20$ παιδιά πόσα λεπτά κοιτούν το κινητό το απόγευμα:\par
\noindent\mbox{}\hfill 35 \ 95 \ 15 \ 90 \ 45 \ 85 \ 40 \ 80 \ 45 \ 50 \ 60 \ 70 \ 55 \ 60 \ 65 \ 25 \ 30 \ 55 \ 35 \ 75\hfill\mbox{}\par
Ομαδοποίησε σε $5$ κλάσεις πλάτους $20$ ξεκινώντας από το $0$ και φτιάξε πίνακα συχνοτήτων και σχετικών συχνοτήτων.

\ex \textbf{Οι βαθμοί.} Η κυρία Αλεξάνδρα έβαλε βαθμούς (στα $20$) σε $25$ μαθητές:\par
\noindent\mbox{}\hfill 10 \ 13 \ 15 \ 11 \ 19 \ 8 \ 14 \ 6 \ 9 \ 18 \ 15 \ 9 \ 12 \ 17 \ 12 \ 16 \ 8 \ 17 \ 12 \ 11 \ 13 \ 13 \ 11 \ 14 \ 7\hfill\mbox{}\par
Φτιάξε πίνακα με κλάσεις πλάτους $4$ από το $4$ ως το $20$. Ποιο ποσοστό των μαθητών πήρε τουλάχιστον $12$;

\ex \textbf{Οι πόντοι.} Ο Νίκος σημείωσε πόσους πόντους έβαλε σε $20$ αγώνες μπάσκετ:\par
\noindent\mbox{}\hfill 10 \ 13 \ 11 \ 8 \ 6 \ 8 \ 21 \ 14 \ 19 \ 2 \ 4 \ 12 \ 12 \ 9 \ 7 \ 18 \ 5 \ 15 \ 16 \ 23\hfill\mbox{}\par
Ομαδοποίησε σε κλάσεις πλάτους $5$ από το $0$. Σε πόσους αγώνες έβαλε λιγότερους από $10$ πόντους και ποιο ποσοστό είναι αυτό;

\ex \textbf{Η βιβλιοθήκη.} Σε μια βιβλιοθήκη γειτονιάς γράφτηκαν $80$ παιδιά. Ο πίνακας έχει σβηστεί εν μέρει. Συμπλήρωσέ τον.
\begin{center}
\renewcommand{\arraystretch}{1.2}
\begin{tabular}{ccc}
\toprule
\textbf{Ηλικία (έτη)} & \textbf{Συχνότητα} & \textbf{Σχετική συχνότητα}\\
\midrule
$6-8$ & $12$ & $\ldots$\\
$8-10$ & $\ldots$ & $25\%$\\
$10-12$ & $32$ & $\ldots$\\
$12-14$ & $\ldots$ & $20\%$\\
\midrule
\textbf{Σύνολο} & $80$ & $\ldots$\\
\bottomrule
\end{tabular}
\end{center}

\ex \textbf{Οι σχολικές τσάντες.} Το βάρος της τσάντας $40$ παιδιών, σε κιλά, δίνει συχνότητες $6$, $12$, $14$, $6$, $2$ στις κλάσεις $2-4$, $4-6$, $6-8$, $8-10$, $10-12$. α) Πόσα παιδιά έχουν τσάντα κάτω από $6$ kg και ποιο ποσοστό είναι; β) Ποιο ποσοστό έχει τσάντα τουλάχιστον $8$ kg; γ) Πόσα παιδιά έχουν τσάντα από $4$ ως $8$ kg (χωρίς το $8$);

\zoni{Το σκέφτομαι}{γιατί και πώς}

\ex Ο Στέλιος ομαδοποιεί τα $25$ ύψη σε $12$ κλάσεις πλάτους $2$ cm και η Ρήνη σε $2$ κλάσεις πλάτους $15$ cm. Τι πρόβλημα έχει κάθε επιλογή; Ποιος αριθμός κλάσεων είναι προτιμότερος;

\ex Ο Βασίλης έφτιαξε τον πίνακα για $20$ παρατηρήσεις: κλάσεις $0-10$, $10-20$, $20-30$, $30-40$ με συχνότητες $5$, $8$, $6$, $1$ και σχετικές συχνότητες $25\%$, $40\%$, $30\%$, $10\%$. Βρες το λάθος και διόρθωσέ το.

\ex Βρες $10$ τιμές για τον χρόνο ενός παιδιού στην ουρά του κυλικείου (σε λεπτά, από $0$ έως $9$) ώστε στον πίνακα με κλάσεις $0-3$, $3-6$, $6-9$ οι συχνότητες να βγουν $2$, $5$, $3$. Υπάρχει μόνο μία λύση;

\ex Η κλάση $160-165$ έχει συχνότητα $4$. Γράψε δύο διαφορετικά πιθανά σύνολα ύψων που δίνουν αυτή τη συχνότητα. Τι δεν μπορεί να μας πει ο πίνακας;

\zoni{Για το σπίτι}{μετρώ και ομαδοποιώ}

\ex Μέτρησε το μήκος του πέλματος $10$ ατόμων της οικογένειας ή των φίλων σου, σε εκατοστά, και ομαδοποίησε τις τιμές σε κλάσεις πλάτους $2$ cm. Φτιάξε πίνακα συχνοτήτων και σχετικών συχνοτήτων.

\ex Σημείωσε για μια εβδομάδα πόσα λεπτά διαβάζεις ή μελετάς κάθε μέρα. Ομαδοποίησε τις $7$ τιμές σε $3$ κλάσεις ίσου πλάτους και φτιάξε τον πίνακα.
